% Spanish CV body (Hoja de Vida)

\begin{center}
  {\Huge \textbf{Jose David Sanabria Aponte}} \\[3pt]
  {\large \textcolor{accent}{\textbf{Desarrollador Backend Senior · Go \& Python · FinTech · Sistemas con IA}}} \\[4pt]
  \small
  +57 317 503 1582 \;$|$\;
  \href{mailto:josedavidsanabriaaponte@gmail.com}{josedavidsanabriaaponte@gmail.com} \;$|$\;
  \href{https://linkedin.com/in/jdsanabriaa}{linkedin.com/in/jdsanabriaa} \;$|$\;
  \href{https://github.com/jdsanabriaap}{github.com/jdsanabriaap} \;$|$\;
  \href{https://jdsa.dev}{jdsa.dev} \;$|$\;
  Bogotá, Colombia
\end{center}

\vspace{-4pt}

\section*{Perfil Profesional}

Ingeniero backend con \textbf{6+ años} en sistemas de pago de alto volumen, plataformas de integración
y automatización con IA. En \textbf{Mercado Libre / Mercado Pago}, lideré integraciones con adquirentes
LATAM en Go — interoperabilidad QR config-driven para \textbf{10+ adquirentes} y onboarding
\textbf{\textasciitilde{}80\%} más rápido con feature flags. Por separado, construí un \textbf{MCP de
documentación de integraciones} con \textbf{prompt-skill de arquitectura hexagonal} para \textbf{todos los
equipos de integraciones} (docs web, Markdown y Mermaid sincronizadas post-release). Como freelance:
\textbf{Holiday Integrator}, automatizaciones independientes y web estática con CI/CD \textbf{OCI +
Cloudflare}. Proyecto personal: \textbf{VibeWindows}. Cursando \textbf{Especialización en IA (UNIR)}.

\section*{Habilidades Técnicas}

\begin{itemize}
  \item \textbf{Lenguajes:} Go (Golang), Python, C\# (.NET / WPF), TypeScript, SQL
  \item \textbf{Backend y Arquitectura:} Microservicios, APIs REST, Arquitectura Hexagonal, Feature Flags, Integraciones Config-Driven, Flujos Transaccionales
  \item \textbf{Datos e IA:} BigQuery, PostgreSQL, MySQL, Visión Artificial (OpenCV, MediaPipe, ONNX), Web Scraping (Selenium, Playwright)
  \item \textbf{Observabilidad y DevOps:} DataDog, Grafana, Git, GitHub Actions, Docker, OCI Object Storage, Cloudflare CDN, Vercel, Azure (AZ-900), Action1 RMM
  \item \textbf{Herramientas IA:} Cursor AI — prompt-skills, bases de contexto, agentes orientados a arquitectura, automatización de docs con Mermaid
  \item \textbf{Metodologías:} Agile / Scrum, CI/CD, trabajo remoto
\end{itemize}

\section*{Experiencia Profesional}

\jobtitle{Desarrollador de Software — Integraciones de Pago}{Mercado Libre / Mercado Pago}{Feb 2024 – Ene 2026}{Remoto · Bogotá}

\begin{itemize}
  \item \textbf{Plataforma de Interoperabilidad QR — Arquitectura Config-Driven:}
        Refactoricé integraciones QR hacia un modelo config-driven (secrets, URLs, adaptadores por contrato),
        centralizando lógica transaccional común — onboarding robusto en \textbf{10+ integraciones LATAM activas}.

  \item \textbf{Arquitectura de Feature Flags — \textasciitilde{}80\% de reducción en tiempos de integración:}
        Feature flags para \textbf{desarrollo en paralelo y testing aislado} de tokenización en 10+
        adquirentes — desacoplando releases del núcleo de los tiempos de integración.

  \item \textbf{Integration Docs MCP — transversal (desarrollo aparte del QR):}
        Construí y publiqué un \textbf{servidor MCP compartido} con \textbf{prompt-skill de arquitectura
        hexagonal} para \textbf{todos los equipos de integraciones de Mercado Pago} (no solo QR).
        Automatiza documentación en \textbf{web y Markdown} con \textbf{diagramas Mermaid} — contexto vivo
        para \textbf{IA y desarrolladores}, actualizado tras merge de branch o release del equipo.

  \item \textbf{Integración Directa Fiserv (sin QR):}
        Actualización de contratos de mensajería MP–Fiserv; validación transaccional extremo a extremo.

  \item \textbf{Diagnóstico Transaccional con BigQuery:}
        Playbooks de diagnóstico reutilizables para inconsistencias MP vs.\ adquirentes.

  \item \textbf{Observabilidad en Producción — DataDog y Grafana:}
        Ampliación de monitoreo; \textbf{releases sin regresiones} por cambio de adquirente.
\end{itemize}

\vspace{4pt}

\jobtitle{Desarrollador de IA y Automatización (Contratista)}{Holiday Experiences S.A.S.}{2020 – Actualidad}{Bogotá}

\textit{Proyectos independientes para el mismo cliente — no un único producto monolítico.}

\begin{itemize}
  \item \textbf{Holiday Integrator} (\texttt{HolidayScraperApp}) — \textbf{C\# / WPF / Playwright:}
        App de escritorio para scraping paralelo en portales B2B (BedsOnline, HotelDo, TravelC, etc.),
        agregación Union-Find, UI mosaico/lista/mapa y handoff de reserva con browser visible.

  \item \textbf{Cumplimiento de Llamadas IFX y Zadarma — Python:}
        Pipeline que eliminó horas de gestión manual de grabaciones — cumplimiento legal sin intervención humana.

  \item \textbf{Sync Transfronteriza de Contratos (Colombia \textrightarrow{} EEUU) — Selenium + OpenCV:}
        Extrae contratos localmente y registra en plataforma US con automatización visual sobre UI legacy.

  \item \textbf{IT Empresarial y RMM:}
        Despliegue y administración de \textbf{Action1 RMM} para parcheo centralizado y seguridad de endpoints.
\end{itemize}

\vspace{4pt}

\jobtitle{Desarrollador Web (Contratista)}{Club House Petlandia}{Ene 2025}{Remoto}

\begin{itemize}
  \item Landing estática multisección (HTML/CSS/JS) — conversión por WhatsApp y SEO local (Madrid, Cundinamarca).
  \item \textbf{CI/CD:} GitHub Actions \textrightarrow{} OCI Object Storage \textrightarrow{} purge Cloudflare en push a \texttt{master}.
  \item Producción: \href{https://petlandiaclubhouse.com}{petlandiaclubhouse.com}
\end{itemize}

\section*{Proyectos Personales}

\begin{itemize}
  \item \textbf{VibeWindows} (v0.2.2) — WPF/.NET 8: control de cursor por gestos; CaptureBroker \textrightarrow{}
        VisionSession \textrightarrow{} motores ONNX o Python TCP; escritorio virtual multi-monitor;
        inyección desactivada por defecto. MVP \textasciitilde{}95\% completado.
\end{itemize}

\section*{Herramientas IA y Prácticas de Ingeniería}

\textbf{Más de 1 año} de \textbf{Cursor AI} en producción — desde prompts básicos hasta \textbf{prompt-skills},
bases de contexto estructuradas y pipelines de documentación con arquitectura hexagonal (web, Markdown,
Mermaid). Aplicado a onboarding de integraciones, servicios Go concurrentes y diagnóstico de datos a escala.

\section*{Educación}

\noindent\textbf{\color{dark}Especialización en Inteligencia Artificial} \hfill \textit{\color{light}Nov 2025 – Actualidad} \\
\textbf{\color{accent}Universidad Internacional de La Rioja (UNIR)} \hfill \textcolor{light}{Remoto}

\vspace{5pt}

\noindent\textbf{\color{dark}Ingeniería de Sistemas} \hfill \textit{\color{light}2017 – 2024} \\
\textbf{\color{accent}Universidad Distrital Francisco José de Caldas} \hfill \textcolor{light}{Bogotá}

\section*{Certificaciones}

\begin{itemize}
  \item \textbf{AZ-900} — Microsoft Certified: Azure Fundamentals
  \item \textbf{SFPC} — Scrum Foundation Professional Certificate
\end{itemize}
